\section{Joint endpoint and incidence localization}
\label{sec:v81-joint}
The critical variable $g=R(-p_0,p_m)$ involves both endpoints.
A one-endpoint estimate loses this geometry. We retain both small
momenta in the characteristic pairing and obtain a gain greater
than one in their common width, at the same time as a superlinear
incidence gain. The two chronological cases require different
estimates; neither reverses the occupation record.

For $0<b<b_0$ let $P_b=M_{\one_{\{|p|<b\}}}$ and retain
$M_s=M_{\one_{\{c<s\}}}$. Take $b_0,s_0$ small and fixed.
All norms below are the density-relative-to-$\nu$ norms of
Section~\ref{sec:v80-incidence}, with stable length exponent $1/8$
and unstable exponent $\gamma_0\le1/100$. In particular the stable
tests carry no additional incidence density; compare
\cite[Section 3.1, equations (10)--(12)]{DPZ}.

\begin{lemma}[A small vector, not a small multiplier]
\label{lem:v81-endpoint-vector}
Uniformly in $b$,
\begin{align}
 \|P_b\one\|&\le Cb^{7/8},& |P_b\one|_w&\le Cb,
                                             \label{eq:v81-small-vector}\\
 \|P_bh\|&\le C\|h\|,& |P_bh|_w&\le Cb^{1/8}\|h\|,
                                             \label{eq:v81-endpoint-multiplier}\\
 |P_bh|_w&\le C|h|_w,& |\ell(P_bh)|&\le Cb^{1/8}\|h\|.
                                             \label{eq:v81-terminal-small}
\end{align}
The small strong norm in \eqref{eq:v81-small-vector} is asserted
only for the particular vector $P_b\one$.
\end{lemma}
\begin{proof}
Near $p=0$, $p=\sin\varphi$ has uniformly bounded derivative and
inverse derivative. A stable graph crosses $|p|<b$ in at most one
interval of length $Cb$. Its strong stable integral is bounded by
$C\min\{|W|,b\}|W|^{-1/8}\le Cb^{7/8}$.
For two matched graphs at distance $d$, the horizontal cuts move by
at most $Cd$. The tests agree in graph coordinates on the common
interval. The difference from unmatched intervals and arclength
Jacobians is at most $C\min(b,d)$, using the separate bound $Cb$
when $d\ge b$. Hence the unstable seminorm is at most
$C\sup_d d^{-\gamma_0}\min(b,d)=Cb^{1-\gamma_0}$.
This proves the strong vector estimate; the weak estimate is the
unweighted intersection-length bound.

The two transverse horizontal cuts satisfy the uniform piecewise
multiplier criterion \cite[Lemma 3.3(b)]{DPZ}. Restricting a weak
test to its retained interval and using the strong stable length
normalization gives the second estimate in
\eqref{eq:v81-endpoint-multiplier}. The same restriction is an
admissible weak test on a shorter stable curve, so multiplication
is also weakly bounded. Mass evaluation is weakly continuous, giving
the last estimate. Smooth averages of the level indicators give
the stated elements and operations in the completions. For the
vector convergence, the same calculation on bands of width tending
to zero gives convergence in the strong norm. No derivative of order
$b^{-1}$ is assigned to the limiting multiplier.
\end{proof}

\subsection{Peak vectors with a physical reference}
Refine the fixed peak charts, without changing the arithmetic
reference, so that
\begin{equation}\label{eq:v81-peak-refinement}
 \frac{\log(1/r_*)}{\log(1/\sigma)}<\frac1{256},\qquad
 \frac{\log(1/r_*)}{\log(1/\rho)}<\frac1{256}.
\end{equation}
Here $\sigma,\rho$ are the constants in
\eqref{eq:v80-spectral-input}. The removed compact annuli join the
fixed-band spectral complement. All unit-modulus points and their
continued peaks remain represented.

On a chart write $\Pi h=v_\xi\ell_\xi(h)$, with uniformly bounded
vectors and functionals. There is no assumption that $\Pi\one$ is
nonzero. Indeed a nontrivial arithmetic peak can annihilate that
particular initial vector.

\begin{lemma}[Localized peak vector and endpoint tests]
\label{lem:v81-peaks}
Put $\theta=31/32$. The chart normalization can be chosen so that
\begin{align}
 \|\Pi h\|+\sup_{d\ge0}|\ell(P_bQ^dh)|
       &\le C\|h\|^{1-\theta}|h|_w^\theta,
                                              \label{eq:v81-interpolation}\\
 |P_bv_\xi|_w&\le Cb^{63/64},&
 |M_sv_\xi|_w&\le Cs^{47/32},&\|M_sv_\xi\|&\le C.
                                              \label{eq:v81-eigen-local}
\end{align}
The constants are independent of the two widths and the length $d$.
\end{lemma}
\begin{proof}
The interpolation proof of Lemma~\ref{lem:v80-interpolation} applies
with the stricter refinement \eqref{eq:v81-peak-refinement} and
exponent $\theta$. In the short-power case use
$|\ell(P_bQ^dh)|\le C|Q^dh|_w\le C(1+d)|h|_w$;
in the long-power case apply the strong inequality at the balancing
time and use power boundedness afterwards. The factor $1+|\log t|$
is absorbed between the exponents $255/256$ and $31/32$.

For completeness we specify the physical vector used to localize
$v_\xi$. At a unit-modulus point the inherited physical peripheral
 theorem, Supplement Theorem~\ref{supp-thm:v40-full-peripheral}, gives
a strong-space eigenvector $u_*$ with $|u_*|=1$ almost everywhere.
Choose a weakly continuous smooth test whose pairing with $u_*$ is
nonzero. Weak continuity of the continued projection makes
$\Pi_\xi u_*$ nonzero on a smaller chart, with that pairing bounded
away from zero. Normalize by this pairing. A finite cover supplies
common constants. The invariant coordinate cylinder is the same
for the circular family, so $u_*$ remains a physical bounded function
when used at neighboring parameters. For all real frequencies,
$|Q^Lu_*|=1$ in the physical realization.
The weak estimates are justified by bounded strong approximants,
not by assigning arbitrary stable-curve traces to an $L^\infty$
representative. The physical Cesaro construction in the cited
peripheral theorem realizes $u_*$ as a strong limit of uniformly
bounded finite Cesaro sums of powers at the reference parameter
applied to one fixed smooth bounded vector, followed by a fixed
normalization. At neighboring real parameters, applying a fixed
power $Q^L$ to these finite physical sums preserves their uniform
modulus bound. Endpoint and homogeneous-incidence intersections
therefore satisfy the same length bounds on every approximant.
Boundedness of the fixed power and multipliers on the common strong
space permits passage to the weak limit. Only after that passage
is the integer $L$ chosen below. Thus the length estimates do not
require a trace assertion for arbitrary bounded measurable densities.

Use
\[
 \Pi_\xi u_*=\lambda^{-L}Q^Lu_*-\lambda^{-L}N^Lu_*.
\]
For $P_b$ choose $L=\lceil\log(1/b)/\log(1/\rho)\rceil$.
The weak norm is at most $Cr_*^{-L}(b+b^{1/8}\rho^L)$,
which is bounded by $Cb^{63/64}$ under
\eqref{eq:v81-peak-refinement}. For $M_s$ choose instead
$L=\lceil(3/2)\log(1/s)/\log(1/\rho)\rceil$.
The homogeneous length estimate, valid for every physical vector
of modulus at most one, gives
$Cr_*^{-L}(s^{3/2}+s^{1/8}\rho^L)\le Cs^{47/32}$.
The bounded strong multiplier gives the last assertion. This
argument neither assigns physical modulus one to a continued
nonperipheral eigenvector nor divides by $\Pi\one$.
\end{proof}

Set
\begin{equation}\label{eq:v81-exponents}
              \alpha=\frac{129}{128},\qquad
              \beta=\frac{17}{16},\qquad
              \delta=\alpha-1=\frac1{128}.
\end{equation}
Let $D_b=\{|p_0|<b,\ |p_m|<b\}$. In the next theorem the endpoint
marks are retained even when the initial and terminal section tests
are dropped in a positive upper comparison.

\begin{theorem}[Two-endpoint, one-incidence local bound]
\label{thm:v81-joint-local}
For every fixed auxiliary band $B_*$ there are $C_{B_*}<\infty$ and
$\rho_{B_*}<1$ such that, uniformly in $R,m\ge2$, $0\le j\le m$,
exact labels, intervals $J$, and $0<b<b_0$, $0<s<s_0$,
\begin{align}
 &m^2\nu\{K_m^{\rm c}=k,A_m=l,L_m\in J,\ D_b,\ c_j<s\}
                                                        \notag\\
 &\quad\le Cb^\alpha s^\beta(|J|+B_*^{-1})
   +C_{B_*}(1+B_*|J|)m^2
                   \min\{b\rho_{B_*}^{m/2},s^2\}.
                                      \label{eq:v81-joint-local}
\end{align}
The leading constant is independent of the previously fixed band.
The estimate is on the original occupation coefficient, not on a
reversed or neighboring occupation class.
\end{theorem}
\begin{proof}
The characteristic pairing is exactly
\[
             \ell(P_bQ^{m-j}M_sQ^jP_b\one).
\]
Suppose first $m-j\ge m/2$. Put $h=M_sQ^jP_b\one$.
Power boundedness and \eqref{eq:v81-small-vector} give
$\|h\|\le Cb^{7/8}$. Its physical realization has modulus at most
one, so $|h|_w\le Cs^{3/2}$. The long-block leading term has
amplitude at most
\[
 Cb^{63/64+(7/8)(1/32)}s^{(3/2)(31/32)}
                  =Cb^{259/256}s^{93/64}.
\]
Its remainder is bounded by $Cb\rho^{m/2}$, using the terminal
factor $b^{1/8}$ as well as the initial strong factor $b^{7/8}$.

If $j\ge m/2$, expand $Q^jP_b\one$ instead. Its projected
coefficient is at most $Cb^{255/256}$ by
\eqref{eq:v81-interpolation}. Set $h=M_sv_\xi$ and $d=m-j$.
There are two bounds for the remaining terminal pairing:
\[
 |\ell(P_bQ^dh)|\le Cb^{1/8},\qquad
 |\ell(P_bQ^dh)|\le Cs^{1457/1024}.
\]
The first uses strong power boundedness and
\eqref{eq:v81-terminal-small}; the second uses
\eqref{eq:v81-interpolation} and \eqref{eq:v81-eigen-local}.
Their geometric interpolation with powers $1/4$ and $3/4$ gives
$Cb^{1/32}s^{4371/4096}$. Thus the leading amplitude in this
case is at most $Cb^{263/256}s^{4371/4096}$.
The remainder is again $Cb\rho^{m/2}$. Both leading bounds imply
$Cb^\alpha s^\beta$, since
$259/256>129/128$ and $4371/4096>17/16$.

The long spectral factor in either case has a four-frequency
Gaussian integral of order $m^{-2}$. On the fixed-band complement
the long block has an exponential gap and the same endpoint norm
factors. Independently, the physical characteristic pairing is at
most $\nu(c_j<s)\le s^2$ at every real frequency. Applying
$|z|\le A+D$, $|z|\le V\Rightarrow |z|\le A+\min(D,V)$
retains the cap in \eqref{eq:v81-joint-local}. The positive interval
majorant and exact torus coefficient extraction now give that
estimate. No number of physical words occurs in the argument.
\end{proof}
